\section{Benchmark Families: 每一类 benchmark 在测什么}

可以把本讲提到的 benchmark 大致归成五类。这样你以后看新榜单时，就知道它在整个地图上的位置。

\begin{table}[H]
\centering
\begin{tabular}{p{2.6cm}p{3.7cm}p{4.1cm}p{3.0cm}}
\toprule
\textbf{家族} & \textbf{测什么} & \textbf{常见代表} & \textbf{常见盲点} \\
\midrule
knowledge & 事实、学科知识 & MMLU, GPQA & 容易被污染和饱和 \\
instruction following & 符合指令 & IFEval, AlpacaEval & judge bias、prompt 敏感 \\
agent & 工具和长流程 & SWE-bench, MLEBench & scaffolding 影响大 \\
reasoning & 抽象推理 & ARC-AGI & 知识和推理难分离 \\
safety & 有害行为与拒答 & HarmBench, AIR-Bench & 情境依赖、地区依赖 \\
\bottomrule
\end{tabular}
\caption{五类 benchmark 的位置感：测的东西不同，失真方式也不同}
\end{table}

\subsection{主要 Benchmark 详细对比}

下面是一张更细致的对比表，包含每个 benchmark 的规模、评估方式和代表性模型得分参考（截至 2025 年中）。注意：分数会随模型版本和评估细节变化，这里只给出数量级和相对关系。

\begin{table}[H]
\centering
\small
\begin{tabular}{p{2.2cm}p{2.0cm}p{1.5cm}p{2.5cm}p{4.5cm}}
\toprule
\textbf{Benchmark} & \textbf{测量能力} & \textbf{题目数} & \textbf{评估方式} & \textbf{代表性分数参考} \\
\midrule
MMLU & 57学科知识 & $\sim$14k & 多选题 log-prob 或生成式 & GPT-4: $\sim$86\%, Llama 3 70B: $\sim$79\% \\
\midrule
MMLU-Pro & 更难知识 & $\sim$12k & 10选项 + CoT & 比 MMLU 低约 15--20 个百分点 \\
\midrule
GPQA & 博士级知识 & $\sim$450 & 多选题 & GPT-4: $\sim$39\%, 人类专家: $\sim$65\% \\
\midrule
IFEval & 指令遵循 & $\sim$500 & 自动约束检查 & 前沿模型: 70--85\% \\
\midrule
Chatbot Arena & 综合偏好 & 持续增长 & 人类 pairwise & ELO 排名（相对分数） \\
\midrule
SWE-bench & 软件工程 & $\sim$2.3k & 单元测试通过率 & 前沿 agent 系统: 30--50\% \\
\midrule
MATH & 数学推理 & $\sim$5k & 精确匹配 & GPT-4: $\sim$52\%, 推理模型: 80+\% \\
\midrule
HumanEval & 代码生成 & 164 & pass@k & GPT-4: $\sim$87\%, 开源模型: 50--75\% \\
\midrule
ARC-AGI & 抽象推理 & $\sim$400 & 精确匹配网格 & 前沿模型 {[}带搜索{]}: $\sim$50\% \\
\midrule
HarmBench & 安全性 & $\sim$500 & 分类器判断 & 因攻击方式差异大，不宜直接比较 \\
\bottomrule
\end{tabular}
\caption{主要 benchmark 详细对比——规模、评估方式和参考分数}
\end{table}

\begin{importantbox}{先判断家族，再判断分数}
如果你连一个 benchmark 属于哪一类都没判断清楚，就直接比较分数，通常只会得到看似精确但实际无意义的结论。
\end{importantbox}

\begin{knowledgebox}{如何把 benchmark 家族和问题匹配起来}
问“这个 benchmark 的失败，是否真的会影响我在乎的结果？”如果答案是肯定的，它才值得作为主指标；否则它最多是辅助信号。
\end{knowledgebox}

